\section{Methods}
\label{sec:methods}

\subsection{Tasks}

We use the prime $p = 113$ throughout, matching \citet{nanda2023progress}.
The three tasks are
\begin{align*}
\textsc{add}: \quad & (a, b) \mapsto (a + b) \bmod p,
                      & a, b &\in \{0, 1, \ldots, p - 1\}, \\
\textsc{sub}: \quad & (a, b) \mapsto (a - b) \bmod p,
                      & a, b &\in \{0, 1, \ldots, p - 1\}, \\
\textsc{mul}: \quad & (a, b) \mapsto (a \cdot b) \bmod p,
                      & a, b &\in \{1, 2, \ldots, p - 1\}.
\end{align*}
For \textsc{mul} we restrict the operands to the multiplicative group
$(\Z/p\Z)^*$. Including zero would cap test accuracy at $1 - 1/p$ for any
otherwise-perfect model, because the network would have to memorise the
$2p - 1$ rows and columns whose product is zero.

\subsection{Tokenisation and input format}

Every example is tokenised as a length-4 sequence
$\bigl[ a, \mathit{op}, b, = \bigr]$. The vocabulary is
$\{0, 1, \ldots, p - 1\} \cup \{+, -, \times, =\}$ for a total of $p + 4$
tokens. The training target is the result digit; the model reads its
output from the last position (the position of $=$). Including the
operation token in the input makes the same network architecture work for
single- and multi-task experiments without modification.

\subsection{Train/test split}

We use a hard 30/70 train/test split per task, sampled by a random
permutation of the $p^2$ (or $(p - 1)^2$) ordered pairs at the configured
seed. For multi-task experiments the split is stratified per task so that
each task contributes 30\% of its own pairs to training. The test set
therefore covers each task uniformly, which lets us track per-task test
accuracy without confounding from class imbalance.

\subsection{Model architecture}

The model is a one-layer decoder-only transformer in the
\citet{nanda2023progress} configuration:
$d_{\mathrm{model}} = 128$, $n_{\mathrm{heads}} = 4$, $d_{\mathrm{head}} = 32$,
$d_{\mathrm{MLP}} = 512$. The embeddings are learned positional and
token embeddings summed at the input. There is no layer normalisation and
no dropout. The MLP is a two-layer feed-forward network with ReLU
activation. The output projection $W_U$ is not tied to the input
embedding $W_E$.

We implement multi-head attention by hand rather than using
\texttt{torch.nn.MultiheadAttention}, so that the per-head residual
contributions are addressable for the ablation studies in Section
\ref{sec:analysis}.

\subsection{Optimiser and training}

We optimise full-batch with AdamW \citep{loshchilov2019adamw}, learning
rate $10^{-3}$, $\beta_1 = 0.9$, $\beta_2 = 0.98$, $\varepsilon = 10^{-8}$,
weight decay $\lambda = 1.0$. Full-batch training is feasible because the
training set has only $\sim$3,800 examples per task at $p = 113,$
$\mathrm{train\_frac} = 0.30$. Step counts vary by experiment:
30,000 for single-task addition and subtraction; 50,000 for single-task
multiplication; 40,000 for multi-task add+subtract; and 60,000 for the
three-task experiment. Three random seeds (42, 137, 271) are run for the
robustness sweep.

\subsection{Metrics}

We log per-step train and test cross-entropy loss and top-1 accuracy.
For multi-task runs these are sliced by task. The headline summary of a
run is the \emph{grok step}: the first step at which test accuracy
exceeds 0.95. This follows the convention of \citet{nanda2023progress}
and lets us compare runs at a glance.

\subsection{Mechanistic-analysis utilities}

\paragraph{Fourier basis.}
Given a trained model we extract the $p \times d_{\mathrm{model}}$
sub-matrix of the embedding consisting of the digit tokens, take the DFT
along the digit axis, and report the per-frequency power
$\sum_d |\widehat W_E(k, d)|^2$. For the multiplicative analysis we first
reorder the rows by powers of a primitive root $g$ of $p$, then take the
DFT. Dominant frequencies are picked greedily until they explain a
configurable fraction (default 90\%) of total non-DC power, with symmetric
pairs $(k, p - k)$ deduplicated.

\paragraph{Feature overlap.}
Two embedding matrices' Fourier feature sets are compared by the cosine
similarity of their per-frequency power vectors. A value of 1 indicates
that the two models put weight on the same frequencies in the same
proportions; 0 indicates orthogonal feature sets.

\paragraph{Head ablation.}
We zero out one attention head's residual contribution at a time and
measure the change in cross-entropy loss. For multi-task models the
ablation is run separately on each task's slice of the data, which
reveals whether heads specialise.

\subsection{Reproducibility}

Every run pins random seeds for both NumPy and PyTorch and writes its
full config (training hyperparameters, model dimensions, dataset
descriptor) as YAML alongside the per-step metrics (NPZ + JSONL). All
experiments run in under an hour on a 2024 Apple-silicon laptop using the
PyTorch MPS backend; no GPU cluster is required.
